\documentclass[platex,dvipdfmx]{jsarticle}

%%%%%%%%%%%%%%%%%%%%%%%%%%%%%%%%%%%%%%%%%%%%%%%%%%%%%%%%%
%						%
%	卒論・修論発表会予稿用Texファイル	%
%   更新：2024/01/17
%						%
%%%%%%%%%%%%%%%%%%%%%%%%%%%%%%%%%%%%%%%%%%%%%%%%%%%%%%%%%


\usepackage{graphicx}
\usepackage{url}

\usepackage{tmu-sd-cs-resume-2024}
\No{A0-00}
\Jtitle{発表会予稿の書き方に関する研究}
\Etitle{A Study of How to Prepare the Manuscript}
\Author{00000000}{東京 太郎}{日野 花子 教授}
\master   % 修士論文発表会の場合
%\bachelor % 特別研究発表会の場合

\begin{document}
\mktitle

\section{はじめに}
この文章は，{\tt tmu-sd-cs-resume-20xx.sty} の使用法を説明しています．以下の説明を読み，フォーマットに関しては間違いのない予稿を作成してください．

\section{予稿執筆環境の準備}

予稿はOverleaf，素の{\LaTeX}またはMicrosoft Wordで執筆できます．

\subsection{Overleafを利用している場合}

予稿はOverleadのテンプレートから利用できます．以下のURLのテンプレートを使うとワンクリックで予稿の執筆環境の準備ができます．論文はmain.tex，参考文献リストはresume.bibに書いてください．

\url{https://www.overleaf.com/latex/templates/tmu-sd-cs-resume/hdydypcwhrdb}

テンプレートを使わないでOverleadで執筆する場合は，空の新規プロジェクトを立ち上げて，別途送信した予稿用スタイルファイル一式をそのプロジェクトにアップロードしてください．latexmkrcファイルが重要です．また，MunuのSettingsのCompilerをLaTexに変更してください．

Overleafでは画面右（プレビュー側）のダウンロードボタンをクリックすることで，PDF化したものがダウンロードされます．これを提出してください．

\subsection{素のTexを利用している場合}

独自の環境でTexを利用している場合も，配布ファイルに含まれる{\tt tmu-sd-cs-resume-20xx.sty}を用いて執筆してください．論文はmain.tex，参考文献リストはresume.bibに書いてください．ファイルはUTF-8でエンコードされています．もしご自身の環境がSHift-JISの場合は変換してください．

\section{予稿の定型記載事項について}

予稿には，フッターに日時・会場等の情報，およびタイトルの先頭に発表番号をつける必要があります．ここでは{\LaTeX}での対応方法を説明します．Microsoft Wordで作成する場合は，同封のTMU\_SD\_CS\_RESUME.docxを使い，元のデザインを崩さないように細心の注意を払って作成してください．記載に必要な情報は{\LaTeX}版と同じですので，以下を参考にしてください．

{\LaTeX}で執筆する場合は，体裁は気にせず，マクロを使って必要な情報を定義してください．具体的にはファイルの先頭部分にある以下の箇所を修正します．．

\begin{quote}
\macrocom{usepackage}{tmu-sd-cs-resume-20xx}\\
\macrocom{No}\{A0-00\}\\
\macrocom{Jtitle}\{発表会予稿の書き方に関する研究\}\\
\macrocom{Etitle}\{A Study of How to Prepare the Manuscript\}\\
\macrocom{Author}\{00000000\}\{首都 太郎\}\{日野 花子 教授\}\\
\macrocom{master}   \% 修士論文発表会の場合\\
\macrocom{bachelor} \% 特別研究発表会の場合
\end{quote}

\macrocom{usepackage}{tmu-sd-cs-resume}はスタイルファイルを指定する部分です．スタイルファイル中に日時や会場名が含まれています．20xxの部分が今年度のファイルになっているかのみ確認してください．

\macrocom{No} は発表番号を設定します．これは当日のプログラムに記載されている情報で，教務委員から通知されます．

\macrocom{Jtitle} および \macrocom{Etitle} は日本語と英語の
論文名をそれぞれ設定します．本文を英語で書く場合でも両方を設定
してください．卒論の場合は英語のタイトルは必要ありません．

\macrocom{Author} は氏名および指導教官を設定します．学習番号，氏名，指導教官の順で記述します．

最後の \macrocom{master} および \macrocom{bachelor} は，どちらか一方のみを残します．修士論文発表会の場合は \macrocom{master} を，特別研究発表会の場合は\macrocom{bachelor} を残します．間違えるとフッターの情報が間違ったものとなってしまうので，注意しましょう．

\section{本文について}
ここでは，本文の記述について説明しています．本文は2ページまで記述することができます（修士論文発表会予稿の場合は4ページも可）．

\subsection{文字サイズなど}
基本的にはスタイルを変更せずに記載する事を前提としますが，長さや読みやすさなどを考慮し，適宜調整する事は可能とします．ただし，周囲の余白は変えないでください．

\subsection{セクションなど}
節などは，\LaTeX のコマンドの \macrocom{section} や \macrocom{subsection} を通常の論文と同様に用いて下さい．

\subsection{図表など}
図や表に関しても特にフォーマットは定めていません．通常の論文と同様に
図\ref{fig:enter-label}のような形式にしてください．

\begin{figure}[tb]
    \centering
    \includegraphics[width=.9\linewidth]{snowy-hino.png}
    \caption{雪の日のキャンパス}
    \label{fig:enter-label}
\end{figure}

表も同様に表\ref{tab:enter-label}のように記載してください．

\begin{table}[htb]
    \centering
    \caption{システムデザイン学部の学科}
    \label{tab:enter-label}
    \begin{tabular}{cl}
        \hline
        \# & 学科名                 \\
        \hline \hline
        1  & 情報科学科              \\
        2  & 電子情報システム工学科    \\
        3  & 機械システム工学科        \\
        4  & 航空宇宙システム工学科     \\
        5  & インダストリアルアート学科  \\
        \hline
    \end{tabular}
\end{table}

\subsection{箇条書き}

予稿はページ数が限られていますので，箇条書き等で端的な説明も有効です．

\begin{itemize}
  \item 一つ目
  \item 二つ目
  \item 三つ目
\end{itemize}

以下のような箇条書きも\LaTeX では簡単に書けます．

\begin{description}
  \item[日時] 一つ目
  \item[場所] 二つ目
  \item[時間] 三つ目
\end{description}

\subsection{数式}

数式は\LaTeX の最も得意とする機能です．論文の文中に$e=mc^2$のように書く事も出来ますし，以下のように独立して書く事も出来ます．

$$ f'(x)=\lim_{x\to h} \frac{f(x+h)-f(x)}{h} $$

$$ \mathrm{e}^{\mathrm{i}\theta} = \cos(\theta) + \mathrm{i}\sin(\theta) $$

$$
A= \left(
\begin{array}{cccc}
a_{11} & a_{12} & \cdots & a_{1n} \\
a_{21} & a_{22} & \cdots & a_{2n} \\
\vdots & \vdots & \quad & \vdots \\
a_{m1} & a_{m2} & \cdots & a_{mn} \\
\end{array}
\right)
$$

\subsection{参考文献}
参考文献はbibtexを使います\cite{codd1970relational}．同封のresume.bibに書き加えるか，自分のbibファイルと差し替えて使ってください．

\section{おわりに}
\LaTeX の書き方に関してはWebにたくさん文書が共有されています．その他，わからない事があれば担当の教員まで問い合わせてください．


\bibliographystyle{acm}
\bibliography{resume}

\end{document}
